\documentclass[11pt]{article}

\usepackage[utf8]{inputenc}
\usepackage[T1]{fontenc}
\usepackage{lmodern}
\usepackage{geometry}
\usepackage{amsmath,amssymb,amsfonts,amsthm,mathtools}
\usepackage{bm}
\usepackage{enumitem}
\usepackage{microtype}
\usepackage{hyperref}
\usepackage{titlesec}
\usepackage{booktabs}
\usepackage{array}
\usepackage{setspace}
\usepackage{mathrsfs}
\usepackage{physics}

\geometry{margin=1in}
\hypersetup{colorlinks=true, linkcolor=black, urlcolor=blue, citecolor=black}
\setlist[enumerate]{topsep=0.4em,itemsep=0.35em}
\setlist[itemize]{topsep=0.3em,itemsep=0.25em}
\titleformat{\section}{\large\bfseries}{\thesection.}{0.5em}{}
\titleformat{\subsection}{\normalsize\bfseries}{\thesubsection.}{0.5em}{}
\setstretch{1.06}

\newcommand{\Aether}{\AE{}ther}
\newcommand{\Aflow}{\AE{}ther-flow}
\newcommand{\TheoryName}{\Aflow{} Interpretation of Relativity}
\newcommand{\TheoryProgram}{\Aether{} / \Aflow{} framework}
\newcommand{\CompletedMetric}{g^{\sharp}}

\newcommand{\ProjectionOperatorCore}{\mathfrak S^{\mathrm{disp,resp,projop}}}
\newcommand{\PullbackBodyResponseCore}{\mathfrak S^{\mathrm{sub,proto,pppresec,pb,bodyresp}}}

\newtheorem{definition}{Definition}
\newtheorem{proposition}{Proposition}
\newtheorem{remark}{Remark}
\newtheorem{corollary}{Corollary}
\newtheorem{theorem}{Theorem}

\title{\textbf{The \TheoryName{}}\\[0.4em]
\large Primitive-Reservoir Observer-Reduction\\
\large Section-Centered Hidden-Response\\
\large Projection-Operator Core\\
\large Next Benchmark Criterion}
\author{}
\date{}

\begin{document}

\maketitle

\begin{abstract}
The current active-primary theorem now proves that the full section-centered displacement-response substrate law still carries benchmark-facing presentational redundancy: once the current-body-realizing substrate body, the displacement projection, the coercive substrate-response operator, the exact-shell discipline, and benchmark faithfulness are fixed, the \(K\)-orthogonal minimal-lift splitting is uniquely forced. The live unexplained stopping point on the recorded section-centered hidden-response route is therefore no longer the full displacement-response substrate law with a separately primitive lift family. It is the smaller current-body-realizing projection-operator core \(\ProjectionOperatorCore\).

This note records the routing consequence of that gain. The next honest benchmark-facing manuscript must now derive or force the current projection-operator core itself, or some nontrivial constituent or relation inside it, from still more primitive explicit substrate structure in a way that changes benchmark-facing status. The claim boundary remains fixed: the one operative metric is still exactly \(\CompletedMetric\), the ordinary matter route is unchanged, there is no second operative metric, the low-energy matter law is not enlarged, and no stronger promotion language is licensed. The positive result is therefore routing discipline below the new projection-operator-core frontier.
\end{abstract}

\section{Introduction}

The active derivational program of \TheoryName{} has again moved the live burden upstream by an honest smaller theorem. The previous current theorem showed that the observer-relevant pullback body-response core \(\PullbackBodyResponseCore\) is forced by a benchmark-faithful section-centered displacement-response substrate law. The new theorem sharpens that result strictly inside the law itself: the separate minimal-lift family is no longer primitive either. What now remains live at current scope is only the smaller current-body-realizing projection-operator core \(\ProjectionOperatorCore\).

That changes what now counts as honest primary progress. Another manuscript that merely carries the same projection-operator core together with its already-forced canonical lift, or that re-expands back to the full displacement-response substrate law as if the lift were still separately primitive, no longer removes any remaining benchmark-facing freedom. The next honest move must therefore target the smaller core itself or some genuinely smaller theorem inside it.

\paragraph{Framework and claim boundary.}
\TheoryProgram{} denotes the broader \Aether{} / \Aflow{} framework built on the claims that the \Aether{} is the underlying four-dimensional substrate of reality, the \Aflow{} is its intrinsic ordered motion, observed three-dimensional space is the local experiential slice of that deeper substrate, S-time is the experienced order of change arising from matter, light, and the \Aflow{}, and observed expansion is the three-dimensional appearance of deeper four-dimensional ordered motion. Within that framework, \TheoryName{} denotes the exact-closure theory in which the effective gravitational dynamics are adopted to be exactly Einsteinian with universal matter coupling. In the active sequence, \emph{adoption} means use of that established relativistic dynamics without claiming substrate derivation, while \emph{derivation} is reserved for a first-principles recovery from explicit substrate variables.

The present note stays strictly inside that boundary. It records only which kind of next manuscript now counts as an honest benchmark-facing gain below the new projection-operator-core frontier.

\section{Fixed Inputs}

\begin{definition}[Projection-operator-core frontier inputs]
The criterion recorded here uses the following fixed inputs.
\begin{enumerate}
    \item The benchmark exact-closure package, which fixes the public one-metric GR target of the repository.
    \item The benchmark gatekeeping layer, including the recorded safeguard that blocks same-output deeper-origin relays from counting as new front-line progress once the live benchmark-relevant datum is unchanged.
    \item The earlier theorem proving that the current observer-relevant pullback body-response core \(\PullbackBodyResponseCore\) is forced by a benchmark-faithful section-centered displacement-response substrate law.
    \item The later theorem proving that the full displacement-response substrate law collapses benchmark-facingly to the smaller current-body-realizing projection-operator core \(\ProjectionOperatorCore\) because the minimal-lift splitting is uniquely forced there.
\end{enumerate}
\end{definition}

\begin{remark}
The present note does not replace those manuscripts. It records the routing consequence of reading them together under the active safeguard against recursive depth drift.
\end{remark}

\section{Why the Live Burden Now Sits Below the Projection-Operator-Core Frontier}

\begin{proposition}[The live burden is renewed below projection-operator-core restatement]
\label{prop:renewed}
At the current stage of the primitive-reservoir observer-reduction line, the projection-operator-core theorem renews the live benchmark-facing burden below any mere restatement of the current-body-realizing projection-operator core \(\ProjectionOperatorCore\).
\end{proposition}

\begin{proof}
The new theorem changes benchmark-facing status because it removes the separate minimal-lift family from the current displacement-response substrate-law frontier as a live unexplained stopping point. Once the lift is forced by \(\ProjectionOperatorCore\), another manuscript that merely preserves the same smaller core no longer removes any remaining benchmark-facing freedom. The live burden therefore no longer sits on restating that same core. It sits on deriving or sharpening the already-active core from still more primitive explicit substrate structure.
\end{proof}

\begin{definition}[Benchmark-neutral same-core relay below the projection-operator-core frontier]
A \emph{benchmark-neutral same-core relay below the projection-operator-core frontier} is a manuscript that preserves the same benchmark one-metric output, the same ordinary matter route, the same current-body-realizing projection-operator core \(\ProjectionOperatorCore\) up to projection-operator-core-faithful equivalence, the same already-forced canonical minimal lift, the same observer-relevant pullback body-response core \(\PullbackBodyResponseCore\), and the same downstream benchmark-facing consequences while merely relocating the unexplained substrate layer deeper, renaming the same smaller core, or re-expanding back to the full displacement-response substrate law or the larger pullback-body-operator, pullback-operator, pullback-quadratic, hidden-response, residual-normal-form, or pre-proto-section presentations without a new derivation, new forcing theorem, or comparably sharp new gain.
\end{definition}

\section{The Current Post-Projection-Operator-Core Criterion}

\begin{definition}[Admissible next benchmark-facing manuscript below the projection-operator-core frontier]
Below the current projection-operator-core frontier, a manuscript counts as an admissible next benchmark-facing move only if it does at least one of the following:
\begin{enumerate}
    \item derives or justifies the current projection-operator core \(\ProjectionOperatorCore\) itself from still more primitive explicit substrate structure in a genuinely benchmark-facing way;
    \item derives or justifies some nontrivial constituent of \(\ProjectionOperatorCore\), such as the substrate body, the displacement projection, the coercive substrate-response operator, or the exact-shell/current-body compatibility relation, from still more primitive explicit substrate structure in a way that sharpens the present route;
    \item proves a sharper necessity, uniqueness, or compatibility theorem for some nontrivial constituent or relation inside \(\ProjectionOperatorCore\);
    \item does one of the previous three while preserving the same benchmark one-metric package, the same ordinary matter route, and the same claim boundary.
\end{enumerate}
\end{definition}

\begin{theorem}[Next honest benchmark-facing task below the projection-operator-core frontier]
The next honest benchmark-facing manuscript below the current projection-operator-core frontier must derive or justify the current smaller core \(\ProjectionOperatorCore\), or some nontrivial constituent or relation inside it, from still more primitive explicit substrate structure in a benchmark-relevant way, or else prove a genuinely sharper theorem inside that same core. A manuscript that only repeats the same current core through another deeper relay, reinstates the already-forced minimal-lift splitting as if it were still primitive, re-expands the discussion back to the full displacement-response substrate law or to the larger pullback-body-operator, pullback-operator, pullback-quadratic, hidden-response, residual-normal-form, or pre-proto-section presentations without such a new gain, or invokes a side branch without doing the same benchmark-facing work does not satisfy the criterion.
\end{theorem}

\begin{proof}
By Proposition \ref{prop:renewed}, the live burden already lies below any mere restatement of the current smaller core \(\ProjectionOperatorCore\). Therefore a subsequent manuscript must improve on that already-active core rather than merely accompany it. A deeper-origin manuscript can do so only if it changes benchmark-facing status by more than depth alone. If it simply reproduces the same core and its same downstream outputs, the routing safeguard classifies it as benchmark neutral. The remaining honest alternatives are therefore exactly those stated in the definition: derive or justify the core itself, derive or justify some nontrivial constituent or relation inside it, or prove a sharper internal theorem there.
\end{proof}

\section{What the Next Move Should Not Be}

\begin{proposition}[Screened next moves below the projection-operator-core frontier]
Below the current projection-operator-core frontier, none of the following is the default next benchmark-facing move:
\begin{enumerate}
    \item another same-output deeper-origin relay that merely preserves the current projection-operator core \(\ProjectionOperatorCore\) up to projection-operator-core-faithful equivalence and therefore merely reproduces the same already-forced canonical lift and the same current body-response core \(\PullbackBodyResponseCore\) together with the same downstream benchmark-facing consequences;
    \item another restatement of the full displacement-response substrate law as if the separately listed minimal-lift family were again independently live;
    \item another re-expansion back to the larger pullback-body-operator, pullback-operator, pullback-quadratic, hidden-response, residual-normal-form, or pre-proto-section presentations as if those already-spent larger presentations were again independently live burdens;
    \item a side branch that does not derive or justify the same projection-operator core, derive or justify a nontrivial constituent or relation inside it, or close a benchmark derivation gate in a genuinely new way;
    \item a manuscript that introduces a second operative metric, enlarges the ordinary matter law, or uses stronger promotion language.
\end{enumerate}
\end{proposition}

\begin{proof}
Item (1) fails the preceding criterion because it leaves the renewed live burden unchanged. Item (2) likewise fails because it merely reintroduces a primitive lift datum that the current theorem already proves to be forced by the smaller core. Item (3) fails because it re-expands to larger already-spent presentations rather than deriving or sharpening the already-active smaller core. Item (4) does not directly address the live burden at current scope unless it performs the same benchmark-facing work named above. Item (5) violates the fixed claim boundary of the benchmark package and therefore cannot count as a disciplined next step on the active line. The benchmark package remains the exact one-metric GR package already fixed by adoption, so any admissible next manuscript must preserve that boundary.
\end{proof}

\begin{remark}
This screening statement is not a denial that those deeper or alternative manuscripts may matter strategically. It is a routing statement: they are not the default way to spend the next benchmark-facing move below the projection-operator-core frontier unless they do the same benchmark-facing work named above.
\end{remark}

\section{Conclusion}

The positive result of this note is routing discipline below the projection-operator-core frontier. The present theorem collapsing the full section-centered displacement-response substrate law to the smaller current-body-realizing projection-operator core \(\ProjectionOperatorCore\) is now the live benchmark-facing gain because it removes the separately primitive minimal-lift splitting as a live unexplained stopping point on the recorded route.

The scope remains narrow and honest. The benchmark package is unchanged: the one operative metric is still exactly \(\CompletedMetric\), the ordinary matter route is unchanged, there is no second operative metric, and the low-energy matter law is not enlarged. Nothing here upgrades adoption to a completed final derivation or licenses stronger promotion language.

The next open burden is therefore clear. The next benchmark-facing manuscript should derive or justify the current projection-operator core itself, or some nontrivial constituent or relation inside it, from still more primitive explicit substrate structure in a way that changes benchmark-facing status. Another same-core deeper relay, another reinstatement of separately primitive lift data, a replay of the full displacement-response substrate law, a re-expansion back to the larger already-spent presentations, or a side continuation that does not do that same benchmark-facing work is not the clean next manuscript target at current scope.

\end{document}
